\section{Introduction}
\label{sec:intro}

Parameter-efficient continual learning (PECL) with LoRA~\citep{hu2022lora} has become a practical
approach for adapting large pretrained vision models to sequential tasks without catastrophic
forgetting~\citep{Wang_2022_CVPR,Gao_2023_ICCV,Smith2023CODA}.
By keeping the backbone frozen and maintaining a single shared low-rank adapter,
SeqLoRA minimizes memory and compute overhead.
The cost is geometric: as new tasks arrive, the adapter drifts away from parameter configurations
that were optimal for old tasks, with no structural anchor preventing this drift.

\textbf{Two families of solutions, one unresolved tension.}
Prior work offers two complementary toolkits for this problem.
The first targets \emph{plasticity}: sharpness-aware minimization (SAM, GAM)
finds flat minima that generalize well within a task~\citep{foretSharpnessAwareMinimizationEfficiently2021,
zhengGradientAwareminimizationDistributed2021} and have been shown to correlate with
reduced forgetting~\citep{NEURIPS2020_518a38cc,shiOvercomingCatastrophicForgetting2021}.
The second targets \emph{stability}: EWC-style Fisher regularization
protects parameter directions important to past tasks~\citep{EWC2017,zheng2025ewclora}.
A natural question is whether combining them is principled or merely heuristic.
Naive combination---applying a sharpness-aware optimizer to the sum
$\Ltask + \Lfisher$---couples the two objectives so that the flatness perturbation
is influenced by the Fisher penalty, obscuring which component is doing the work.

\textbf{The geometric question.}
We reframe this as a question about gradient geometry:
do the AS$^1$ flatness gradient $\ggam - \gclean$ and the past-Fisher gradient $\gfisher$
point in \emph{similar} or \emph{orthogonal} directions in parameter space?
If they are nearly orthogonal, a joint update combines two genuinely independent signals
and each component can be designed and tuned independently.
If they are aligned (or anti-aligned), the two objectives substantially interfere.

\textbf{Our finding.}
We measure the cosine similarity between the AS$^1$ gradient component
and the Fisher gradient component throughout training on six datasets.
The result is unambiguous: $\cosff \approx -0.02$ to $-0.001$ across all datasets,
all tasks, and all $\lambda$ settings (Figure~\ref{fig:cosine_histogram}).
This near-zero cosine is not an accident.
We provide a theoretical explanation: when AS$^1$ targets the current-task Hessian's
principal subspace and Fisher-EWC targets the past-task Fisher's principal subspace,
the spectral separation between current and past tasks drives the cosine toward zero.

\begin{figure}[t]
\centering
\fbox{\rule[-.5cm]{0cm}{4cm} \rule[-.5cm]{8cm}{0cm}}
\caption{Distribution of $\cosff$ (cosine between AS$^1$ flatness gradient and Fisher gradient)
across all tasks and six datasets. Values cluster near zero, confirming that the two components
address orthogonal subspaces of the gradient space. This enables a principled joint design
that treats them as independent signals.}
\label{fig:cosine_histogram}
\end{figure}

\textbf{Method: FS-LoRA.}
Building on this orthogonality observation, we propose \textbf{Flat-Stable LoRA (FS-LoRA)},
a method with three design principles:
\begin{enumerate}
  \item \textbf{Adapter-space Fisher-EWC.}
    We define the Fisher penalty and the drift reference point in the effective update space
    $\DeltaW = BA$, not in the raw factor coordinates $(A, B)$.
    This makes the regularization basis-invariant with respect to the LoRA factorization.
  \item \textbf{Decoupled AS$^1$ and Fisher gradients.}
    The sharpness-aware optimizer (GAM) is applied only to $\Ltask$, not to
    $\Ltask + \Lfisher$.
    The Fisher gradient is added as an independent term.
    This preserves the orthogonality structure and makes each component interpretable.
  \item \textbf{Trace-normalized Fisher.}
    Raw Fisher values in $\DeltaW$-space are typically $\mathcal{O}(10^{-6})$.
    We normalize by the trace to make $\lambdaewc$ a dimensionless relative weight,
    ensuring the Fisher penalty is consistently active across datasets.
\end{enumerate}

The resulting update rule is:
\begin{equation}
  \gupdate = \gclean + \lambdaflat(\ggam - \gclean) + \gfisher.
  \label{eq:update_rule}
\end{equation}

\textbf{What this paper is, and is not.}
FS-LoRA is not simply ``C-Flat + EWC-LoRA stacked together.''
C-Flat~\citep{NEURIPS2024_C-Flat} adds flatness control to CL methods but does not address Fisher
regularization and does not study the gradient interaction between flatness and stability.
EWC-LoRA~\citep{zheng2025ewclora} provides the stability component but has no flatness control.
The critical question---\emph{whether combining them is principled or heuristic}---requires knowing
whether the two gradient signals interfere.
The orthogonality finding ($\cosff \approx 0$) is the answer to that question,
and it is the primary contribution of this paper.
Prior flatness-plus-orthogonality work~\citep{yangRevisitingFlatnessAwareOptimization2025,garg2026fopng}
addresses orthogonality between \emph{task-task} gradients, enforced as a constraint;
FS-LoRA discovers orthogonality between the \emph{flatness} and \emph{Fisher} gradient components
as an emergent geometric property.

\textbf{Contributions.}
\begin{enumerate}
  \item \textbf{Empirical discovery of near-orthogonality (primary contribution).}
    We provide the first systematic measurement of the cosine between the AS$^1$ flatness gradient
    component $(\ggam - \gclean)$ and the Fisher-EWC gradient $\gfisher$
    in LoRA-based continual learning, confirming $|\cosff| \leq 0.03$ across
    six datasets, all tasks, and all $\lambdaewc$ settings.
    This near-zero cosine is not imposed by design: it emerges from the spectral
    structure of the current-task Hessian and past-task Fisher in adapter space.
  \item \textbf{Geometric explanation.}
    We characterize orthogonality via the spectral gap between $H_{\DeltaW}$ and $F_{t'}$
    in the adapter update space (Proposition~\ref{prop:orthogonality}), and show that
    coupling the SAM closure with $\Lfisher$ destroys this structure by injecting
    a $J^\top F_{t'} J\epsilon$ cross-term.
  \item \textbf{FS-LoRA method.}
    Building on the orthogonality finding, we propose a principled joint optimizer with
    decoupled gradients, $\DeltaW$-space Fisher, and trace normalization.
    Each design choice is motivated by the geometry, not by ablation search.
  \item \textbf{Multi-dataset evaluation.}
    Experiments on six fine-grained vision datasets confirm consistent improvements
    in final average accuracy and forgetting relative to SeqLoRA+GAM.
    FS-LoRA matches or exceeds EWC-LoRA on standard PECL benchmarks
    (CIFAR-100, ImageNet-R, DomainNet, ImageNet-A) with LR-tuned settings.
\end{enumerate}
